{\scriptsize
    \begin{minipage}[t]{0.5\textwidth}
        \begin{tabularx}{\linewidth}{|p{0.53\textwidth}|X|}
            \hline 
            \textbf{Notions et contenus} & \textbf{Capacités exigibles} \\
            \hline 
            6.4. Introduction à la physique du laser & \\
            \hline 
            6.1. Milleu amplificateur de lumière & \\
            \hline 
            Absorption, émission stimulée, émission spontanée. & Distinguer les propriétés d'un photon émis par émission spontanée ou stimulée. \\
            \hline 
            Coefficients d'Einstein. & Associer l'émission spontanée à la durée de vie d'un niveau excité. Utiliser les coefficients d'Einstein dans le cas d'un système à plusieurs niveaux non dégénérés. \\
            \hline 
            Amplification d'ondes lumineuses par émission stimulée. & Justifier qualitativement la nécessité d'une inversion de population pour parvenir à amplifier une onde électromagnétique dans un laser. \\
            \hline
        \end{tabularx}
\end{minipage}%
\begin{minipage}[t]{0.5\textwidth}
    \begin{tabularx}{\linewidth}{|X|X|}
        \hline 
        \multicolumn{2}{|l|}{Propriėtés optiques d'un faisceau spatialement limité } \\
        \hline 
        Description simplifiée d'un faisceau de profil gaussien : waist, longueur de Rayleigh, ouverture angulaire. & Justifier qualitativement l'inadéquation du modèle de l'onde plane pour décrire un faisceau laser. Utiliser l'expression fournie du profil radial d'intensité. Construire l'allure d'un faisceau de profil gaussien à partir de l'enveloppe d'un faisceau cylindrique et d'un faisceau conique. Exploiter qualitativement le phénomène de diffraction pour relier le waist et l'ouverture angulaire du faisceau à grande distance. \\
        \hline 
        Transformation à l'aide d'une lentille d'un faisceau cylindrique en faisceau conique et réciproquement. Élargisseur de faisceau. & Déterminer la dimension et la position de la section minimale du faisceau émergeant d'une lentille éclairée par un faisceau cylindrique. \\
        \hline
    \end{tabularx}
\end{minipage}
}
